\documentclass[sigconf]{acmart}

% ACM acmart sigconf template. Same template family as SIGIR and
% CIKM (both use acmart sigconf); confirm conference-specific
% \acmConference{} fields before camera-ready.

\usepackage{booktabs}
\usepackage{multirow}
\usepackage{amsmath, amssymb}
\usepackage{xcolor}
\usepackage{enumitem}
\usepackage{url}

\setcopyright{none}
\copyrightyear{2026}
\acmYear{2026}
% TODO before camera-ready: replace with the venue-specific
% \acmConference{} fields. Current placeholder is sufficient for
% submission; the publisher provides the exact strings on acceptance.
\acmConference[Reproducibility Track]{Reproducibility Track Submission}{2026}{}
\acmBooktitle{Reproducibility Track Submission, 2026}

\begin{document}

\title{Reproducing and Auditing GCCP: Undocumented Hazards and Unexpected Sensitivities in Anchor-Based Pointwise Ranking}

\author{Utshab Kumar Ghosh}
\affiliation{%
  \institution{Missouri University of Science and Technology}
  \city{Rolla}
  \state{MO}
  \country{USA}
}
\orcid{0000-0003-3096-6909}
\email{u.ghosh@mst.edu}

\author{Shubham Chatterjee}
\affiliation{%
  \institution{Missouri University of Science and Technology}
  \city{Rolla}
  \state{MO}
  \country{USA}
}
\orcid{0000-0002-6729-1346}
\email{shubham.chatterjee@mst.edu}

\renewcommand{\shortauthors}{Ghosh and Chatterjee}

\begin{abstract}
We reproduce and audit GCCP/PAGC (Long et al., SIGIR 2025), an
anchor-based pointwise LLM reranker. Our headline result is that
\textbf{a single undocumented choice -- the case of the
yes/no target tokens -- collapses NDCG@10 from 0.66 to 0.24} in
isolation on DL19/Flan-T5-Large; seven other equally fail-silent
implementation details together explain the rest of the gap between
a paper-only reimplementation and a working pipeline. After
recovering all eight, we reproduce the paper's
2-component PAGC = RG-YN+GCCP within $1.6\%$ mean absolute NDCG@10
on TREC-DL and $1.9$--$4.5\%$ on the 8 paper-aligned BEIR sets at
three Flan-T5 scales, exceeding the paper on five
(\textsection\ref{sec:reproduction}).
We extend the original paper's per-cell paired $t$-tests (daggers in
its Tables~1--2, ``$+$\,GCCP vs single component'' only) with paired
bootstrap (10{,}000 resamples) + Holm-Bonferroni correction across
three comparison families and 22 primary
(dataset $\times$ retrieval $\times$ model) settings. PAGC vs RG-YN
remains Holm-significant in 12/22 settings (all positive) -- the
paper's headline ordering holds. \textbf{What does \emph{not} hold
under correction is PAGC-vs-GCCP-alone}: Holm-significant in only
5/22 settings, with mixed signs, including a Holm-significant
\emph{negative} on the largest BEIR set (DBPedia-Entity, $n_q{=}400$,
$\Delta{=}{-}0.0144$, $p_{\text{Holm}}{=}0.032$). Aggregation with
RG-YN is unambiguously useful only under BM25; under E5 dense
retrieval it is roughly a wash and on at least one large-$n_q$
benchmark it actively hurts (\textsection\ref{sec:stats},
\textsection\ref{sec:beir-e5}). Three further findings of
independent interest: (a) the spectral MDS anchor that the paper
recommends never wins on any of 8 BEIR sets we tested and finishes
\emph{last} on 3/8; (b) replacing BM25 with off-the-shelf E5
retrieval already exceeds the paper's reported PAGC-QSG (3-component)
numbers on TREC-DL with Flan-T5-XL, and the contrastive anchor
contributes a smaller marginal gain when the candidate list is
already clean; (c) the contrastive-anchor mechanism transfers to a
4-bit AWQ-quantized 72B decoder-only model on a single 48~GB GPU,
reaching PAGC NDCG@10 $=0.7465$ on DL19, exceeding the paper's
reported Flan-UL2 (both the RG-YN+GCCP and PAGC-QSG variants); on
a 50-query BEIR check, this 72B backbone reproduces the strict
PAGC$>$GCCP$>$RG-YN ordering on SciFact and \emph{strengthens} the
DBPedia-Entity negative finding -- here PAGC underperforms not
just GCCP-alone but the E5 first-stage retriever itself, so the
negative result survives a complete change of backbone family,
scale, and precision.
\end{abstract}

%% CCS concepts (ACM 2012)
\begin{CCSXML}
<ccs2012>
   <concept>
       <concept_id>10002951.10003317.10003347</concept_id>
       <concept_desc>Information systems~Retrieval models and ranking</concept_desc>
       <concept_significance>500</concept_significance>
   </concept>
   <concept>
       <concept_id>10010147.10010178.10010179</concept_id>
       <concept_desc>Computing methodologies~Natural language processing</concept_desc>
       <concept_significance>300</concept_significance>
   </concept>
   <concept>
       <concept_id>10002951.10003317.10003338.10003344</concept_id>
       <concept_desc>Information systems~Test collections</concept_desc>
       <concept_significance>300</concept_significance>
   </concept>
</ccs2012>
\end{CCSXML}

\ccsdesc[500]{Information systems~Retrieval models and ranking}
\ccsdesc[300]{Computing methodologies~Natural language processing}
\ccsdesc[300]{Information systems~Test collections}

\keywords{LLM ranking, reproducibility, zero-shot retrieval, BEIR,
TREC, statistical significance, contrastive prompting}

\maketitle

% ============================================================
\section{Introduction}
% ============================================================

Zero-shot LLM-based document ranking has converged on three
prompt families: \emph{listwise} (rank a list at once), \emph{pairwise}
(compare two documents), and \emph{pointwise} (score one document at a
time). The pointwise family is the cheapest -- $O(n)$ inferences for
$n$ candidates -- but the documents are scored independently, so the
LLM never sees how they relate to each other. Long et al.~\cite{long2025precise}
propose GCCP to bridge this gap: they insert an \emph{anchor} passage
-- a query-focused summary built from the top-$m$ candidates by spectral
multi-document summarization -- into a pointwise prompt and ask the
LLM to compare the candidate against the anchor. Aggregating the
resulting score with a standard relevance-grading score (PAGC)
yields, in their paper, listwise-competitive NDCG@10 at pointwise
cost.

The method is attractive because it is training-free, model-agnostic
in principle, and has no $O(n^2)$ explosion. We selected it for
reproduction because (i) it is recent, well-cited, and proposes a
clean architectural idea, (ii) the authors release code, allowing a
controlled audit of paper-vs.-implementation discrepancies, and (iii)
the method is small enough to be \emph{re-implementable} from the
paper alone in a few weeks of independent work, which is exactly
the regime where under-documented implementation choices matter
most.

This paper makes the following contributions:

\begin{itemize}[leftmargin=1.5em,nosep]
    \item A paper-then-code reimplementation: we coded each of the
          four pipeline components from the paper text first, then
          consulted the authors' release only to resolve ambiguities
          (with each consultation logged in our project plan). The
          result reproduces the paper's headline numbers with a
          mean absolute gap of $1.6\%$ across the six TREC-DL cells
          and $1.9$--$4.5\%$ across the 8 paper-aligned BEIR sets
          at three model scales, exceeding the paper on five cells
          (TREC-DL/Flan-T5-XL/DL19, TREC-DL/Flan-T5-Large/DL20,
          BEIR-SciFact/Flan-T5-Large, and BEIR-Touch\'e-2020 at
          both Flan-T5-XL and Flan-UL2).

    \item A catalogue of eight undocumented implementation details
          discovered during reproduction (Section~\ref{sec:undocumented}),
          including a single-token case discrepancy that drops NDCG@10
          from 0.66 to 0.24.

    \item Paired-bootstrap significance tests with Holm-Bonferroni
          correction across three comparison families and 22 primary
          settings. The original paper reports only per-cell paired
          $t$-tests for ``$+$\,GCCP vs single component'' (daggers in
          its Tables~1--2), without multiple-comparison correction or
          PAGC-vs-GCCP-alone tests. We find PAGC vs RG-YN
          Holm-significant in 12/22 (all positive); GCCP vs RG-YN
          in only 3/22; \textbf{PAGC vs GCCP-alone in 5/22 with
          mixed signs} -- 4 positive (BM25) and 1 \emph{negative}
          (DBPedia-Entity under E5, $\Delta{=}{-}0.0144$,
          $p_{\text{Holm}}{=}0.032$).

    \item Four extensions absent from the original paper:
          (i) E5 dense retrieval as the first stage on TREC-DL and
          all 8 BEIR sets; (ii) decoder-only LLMs adapted via chat
          templates -- LLaMA-3.1-8B-Instruct, Qwen-2.5-7B-Instruct,
          Mistral-7B-Instruct-v0.3, and
          Qwen-2.5-72B-Instruct-AWQ (4-bit, single 48~GB GPU);
          (iii) a 3-component PAGC-RS-YN-GCCP variant adding the
          authors' RG-S$(0,4)$ classifier (Section~\ref{sec:pagc_qsg});
          and (iv) a positioning of our 2-component PAGC against the
          listwise/pairwise baselines from the paper's Table~2 at
          the same Flan-T5 scale (Section~\ref{sec:rankgpt}).

    \item Aggregation and anchor-construction ablations: the
          paper's $\alpha{=}0.5$ linear baseline is statistically
          tied with Borda voting on the 8 cells we test. The
          spectral MDS anchor never wins on any of 8 BEIR sets and
          finishes \emph{last} on 3/8 -- a stronger negative
          finding than the paper's TREC-DL exception.
\end{itemize}

All code, per-query scores, statistical tests, and run logs are
available at \url{https://github.com/utshabkg/GCCP-reproduce}.

% ============================================================
\section{Background: GCCP and PAGC}
% ============================================================

GCCP stands for \textbf{Global-Consistent Comparative Pointwise};
PAGC stands for \textbf{Post-Aggregation with Global Context}.

\paragraph{Notation.} Given a query $q$ and a candidate passage list
$D = \{d_1, \dots, d_n\}$ from BM25 first-stage retrieval, the goal is
to produce a re-ranking $\pi$ of $D$ that maximizes NDCG@10 against
graded relevance judgments.

\paragraph{RG-YN baseline.} The standard pointwise relevance-grading
prompt asks the LLM whether a passage is relevant to the query and
scores it by $f_{\text{RG-YN}}(q,d) = P(\text{`yes'} \mid q,d) /
(P(\text{`yes'} \mid q,d) + P(\text{`no'} \mid q,d))$ at the LLM's
output position (Eq.~3 of \cite{long2025precise}).

\paragraph{Anchor construction.} Long et al.\ build an anchor $d_a$
from the top-$m$ BM25 documents using spectral
multi-document summarization (Eqs.~5--9): sentences are embedded with
TF-IDF, an affinity matrix is thresholded at $\theta$, the normalized
Laplacian is decomposed, and the Fiedler vector partitions sentences
into two clusters; the anchor is the larger cluster trimmed to $z$
sentences, ordered by source document.

\paragraph{GCCP.} For each candidate $d_i$, the LLM is prompted with
``Given a query $q$, which of the following two passages is more
relevant: Passage A or Passage B?'' with $d_i$ in slot A and $d_a$ in
slot B. The contrastive score is
$f_c(q,d_i,d_a) = P(\text{`A'}) / (P(\text{`A'}) + P(\text{`B'}))$
(Eq.~10).

\paragraph{PAGC.} Final score is the equal-weight mean of RG-YN and
GCCP:
$f_{\text{PAGC}}(q,d) = \tfrac{1}{2}\bigl(f_{\text{RG-YN}}(q,d) + f_c(q,d,d_a)\bigr)$
(Eq.~11). The released code additionally min-max normalizes each
component's scores to $[0,1]$ within the candidate list of $q$
before averaging; this normalization step is not in Eq.~11 of the
paper but is required for the components' scales to match. We
catalogue this as undocumented detail \#8 in
Section~\ref{sec:undocumented}.

\paragraph{Hyperparameters.} The paper specifies
$m=10$ (anchor pool), $z=10$ (anchor size). The threshold $\theta$,
BM25 parameters, decoder-input format, and target-token case are
\emph{not} specified -- one of the central findings of this paper.

% ============================================================
\section{Reproduction Methodology}
% ============================================================

\subsection{Implementation}
We re-implemented the four pipeline components (RG-YN, spectral MDS,
GCCP, PAGC) from the paper text alone in the first pass, which gave
NDCG@10 $\approx 0.24$ on DL19 RG-YN against the paper's $0.66$. The
authors' released code\footnote{\url{https://github.com/ChainsawM/GCCP}}
was then consulted to resolve ambiguities (with each consultation
logged in our project plan); each undocumented detail recovered in
this second pass is itemized in Section~\ref{sec:undocumented}. The
resulting modules total $\sim$1.2k lines of Python.

\paragraph{Scope of reproduction.} The original paper introduces
several PAGC variants that combine GCCP with different pointwise
backbones (RG-YN, RG-S(0,4), QG, and combinations thereof). The
\emph{strongest} variant in the original paper's Table~2 is
PAGC\textsubscript{QSG} = QG $+$ RG-S $+$ GCCP. To keep the
inference budget tractable we focus on the 2-component
RG-YN $+$ GCCP variant throughout the main pipeline; we add a
3-component PAGC-RS-YN-GCCP check in
Section~\ref{sec:pagc_qsg}. We do not implement QG -- it requires a
different inference pattern (full-sequence query log-likelihood
given the document) and we leave it for follow-up work.

\subsection{Environment}
Python 3.10, PyTorch 2.1.0+cu121, \texttt{transformers} 4.36 for the
encoder-decoder pipeline, \texttt{transformers} 4.46 for the
decoder-only extension. Hardware: $2\times$NVIDIA RTX 6000 Ada (48~GB).
NDCG@10 computation uses \texttt{pytrec\_eval}~\cite{vangysel2018pytreceval}, which matches the
official \texttt{trec\_eval} binary; we discuss why that choice
matters in Section~\ref{sec:undocumented}.

\subsection{Data}
TREC DL 2019 (43 queries) and DL 2020 (54 queries) over MS MARCO
v1 passages~\cite{bajaj2016msmarco}, retrieved with Pyserini's
prebuilt index using BM25 parameters $k_1{=}0.9$, $b{=}0.4$.
For BEIR, we use Pyserini's prebuilt Lucene indices over the eight
subsets used by Long et al.: TREC-COVID, Touch\'e-2020, DBPedia,
SciFact, Signal1M, TREC-News, Robust04, NFCorpus.

\subsection{Models}
Three encoder-decoder Flan models from the original paper:
Flan-T5-Large (780M), Flan-T5-XL (3B), and Flan-UL2 (20B). Four
decoder-only models in our extension (Section~\ref{sec:decoder}):
LLaMA-3.1-8B-Instruct, Qwen-2.5-7B-Instruct, Mistral-7B-Instruct-v0.3,
and Qwen-2.5-72B-Instruct-AWQ (4-bit AWQ~\cite{lin2024awq}-quantized, on a single
48~GB GPU). The Flan and 7--8B decoder-only models run in FP16; the
72B-AWQ variant runs at 4-bit weight precision with FP16 activations.
Inference is sequential (no batching), matching the original code;
we discuss latency in Section~\ref{sec:efficiency}.

% ============================================================
\section{Reproduction Results}
% ============================================================
\label{sec:reproduction}

\subsection{TREC DL}
Table~\ref{tab:dl-pagc} reports PAGC NDCG@10 on TREC DL 2019/2020
across the three model scales. The mean absolute gap is
$1.6\%$ across the six cells. We exceed the paper on
\textbf{2/6 cells} -- DL19/Flan-T5-XL ($+0.9\%$) and
DL20/Flan-T5-Large ($+3.7\%$) -- and remain within $1\%$ on three
others. The largest negative gap is $-2.5\%$ on
DL19/Flan-T5-Large. Note that ``Paper'' in the table refers to
the RG-YN+GCCP row of Table~1 in \cite{long2025precise}, which
matches our 2-component PAGC = RG-YN $+$ GCCP construction.

\begin{table}[t]
\centering
\small
\begin{tabular}{llccc}
\toprule
\textbf{Dataset} & \textbf{Model} & \textbf{Ours} & \textbf{Paper} & \textbf{Gap} \\
\midrule
DL19 & Flan-T5-Large & 0.6834 & 0.7012 & $-2.5\%$ \\
DL19 & Flan-T5-XL    & \textbf{0.7030} & 0.6969 & $+0.9\%$ \\
DL19 & Flan-UL2      & 0.7095 & 0.7206 & $-1.5\%$ \\
\midrule
DL20 & Flan-T5-Large & \textbf{0.6515} & 0.6281 & $+3.7\%$ \\
DL20 & Flan-T5-XL    & 0.6760 & 0.6810 & $-0.7\%$ \\
DL20 & Flan-UL2      & 0.7009 & 0.7023 & $-0.2\%$ \\
\bottomrule
\end{tabular}
\caption{PAGC NDCG@10 on TREC DL. ``Ours'' values come from the
original full-pipeline reproduction. Ablation tables that require
per-query scores (Tables~\ref{tab:agg}, \ref{tab:rankgpt},
\ref{tab:stats}) draw from a later score-preserving re-run on the
same configuration; numerical differences across runs are within
$\pm 1$~pt NDCG@10 and statistically indistinguishable, attributable
to small differences in the spectral-MDS anchor random seed. The
$\pm 1$~pt re-run noise does not flip any of the TREC-DL/decoder
Holm decisions in Section~\ref{sec:stats} (all relevant $|\Delta|$
are $\geq 0.04$). The one close call is the BEIR-DBPedia-Entity
PAGC vs GCCP cell ($\Delta=-0.0144$, $p_{\text{Holm}}=0.032$),
where the magnitude is comparable to the run-to-run jitter on the
TREC-DL score-preservation re-run; that BEIR-E5 cell, however, was
\emph{not} re-run with score saving (it is a single fresh run from
the BEIR-E5 pipeline), so the $\pm 1$~pt jitter does not directly
apply -- but a future replication should treat this as a finding
to verify rather than to take as settled.}
\label{tab:dl-pagc}
\end{table}

\subsection{BEIR}
Table~\ref{tab:beir-pagc} reports PAGC NDCG@10 across the eight BEIR
subsets and three model scales (paper numbers from Table~1
RG-YN+GCCP rows of \cite{long2025precise}, the row that corresponds
to our 2-component PAGC). At T5-Large the mean absolute gap is
4.5\% (median $4.4\%$, signed mean $-3.5\%$); at T5-XL,
$1.9\%$ ($1.6\%$ median, $-1.4\%$ signed mean); at UL2,
$3.3\%$ ($2.7\%$ median, $-1.3\%$ signed mean). Median and signed
mean are reported alongside the absolute mean because Touch\'e-2020
($+7.9\%$ at UL2) and SciFact ($+4.2\%$ at T5-Large) are both
positive outliers that pull unsigned means upward. We exceed the
paper on SciFact (T5-Large, $+4.2\%$) and Touch\'e-2020 (both
T5-XL, $+2.1\%$, and UL2, $+7.9\%$); the remaining gaps are
negative but not consistent in magnitude across datasets,
suggesting per-dataset preprocessing differences rather than a
single bug.

\begin{table*}[t]
\centering
\small
\begin{tabular}{lcccccc}
\toprule
 & \multicolumn{2}{c}{\textbf{Flan-T5-Large}}
 & \multicolumn{2}{c}{\textbf{Flan-T5-XL}}
 & \multicolumn{2}{c}{\textbf{Flan-UL2}} \\
\cmidrule(lr){2-3}\cmidrule(lr){4-5}\cmidrule(lr){6-7}
\textbf{Dataset} & Ours & Paper & Ours & Paper & Ours & Paper \\
\midrule
SciFact           & \textbf{0.6403} & 0.6145 & 0.6840 & 0.6854 & 0.7114 & 0.7295 \\
NFCorpus          & 0.3620 & 0.3638 & 0.3728 & 0.3756 & 0.3776 & 0.3792 \\
TREC-COVID        & 0.7294 & 0.7641 & 0.7552 & 0.7761 & 0.7495 & 0.7781 \\
TREC-News         & 0.3820 & 0.4112 & 0.4490 & 0.4740 & 0.4563 & 0.4864 \\
Touch\'e-2020     & 0.2650 & 0.2928 & \textbf{0.3078} & 0.3016 & \textbf{0.3161} & 0.2930 \\
DBPedia           & 0.3898 & 0.4181 & 0.4061 & 0.4149 & 0.4138 & 0.4265 \\
Robust04          & 0.4800 & 0.4914 & 0.5279 & 0.5346 & 0.5347 & 0.5467 \\
Signal1M          & 0.2983 & 0.3027 & 0.3204 & 0.3230 & 0.3164 & 0.3186 \\
\midrule
\textbf{Avg.\ abs.\ gap} & \multicolumn{2}{c}{4.5\%}
                         & \multicolumn{2}{c}{1.9\%}
                         & \multicolumn{2}{c}{3.3\%} \\
\bottomrule
\end{tabular}
\caption{BEIR PAGC NDCG@10 across three model scales.}
\label{tab:beir-pagc}
\end{table*}

% ============================================================
\section{Undocumented Implementation Details}
% ============================================================
\label{sec:undocumented}

A faithful first-pass re-implementation from the paper alone produced
NDCG@10 $\approx 0.24$ on DL19 RG-YN (paper: 0.66). We
audited the authors' released code and found eight undocumented
details that together explain the entirety of the gap; their
individual contributions, fixed in isolation, are:

\begin{enumerate}[leftmargin=1.5em,nosep]
    \item \textbf{Decoder input string for T5}: the authors set
          \texttt{decoder\_input\_ids} to a literal \texttt{`<pad> '}
          for RG-YN and \texttt{`<pad> Passage '} for GCCP. With an
          empty decoder input, T5 produces near-uniform probabilities
          for the candidate tokens.

    \item \textbf{Target-token case (RG-YN)}: lowercase
          \texttt{`yes'/`no'}, not \texttt{`Yes'/`No'}. T5's
          tokenizer produces logits for both, but with a prompt
          ending in ``Answer 'yes' or 'no'`` essentially all
          probability mass goes to the lowercase variant. This single
          change moves NDCG@10 on DL19 RG-YN from 0.24 to 0.55.

    \item \textbf{Target-token case (GCCP)}: uppercase
          \texttt{`A'/`B'}, with the prompt ending ``Output Passage A
          or Passage B:''. The capitalization flips here relative to
          item~2.

    \item \textbf{Spectral threshold $\theta = 0.2$}: not stated in
          the paper. With $\theta = 0$ (no thresholding) the affinity
          graph is dense and the Fiedler partition becomes uninformative.

    \item \textbf{BM25 parameters $k_1{=}0.9$, $b{=}0.4$}: not stated.
          Defaulting to Pyserini's stock $k_1{=}0.82$, $b{=}0.68$
          changes the candidate pool composition.

    \item \textbf{Document truncation to 128 tokens}: not stated. Without
          it, the encoder is overwhelmed by a few long candidates and
          the prompt template no longer fits.

    \item \textbf{NDCG implementation}: the paper does not specify
          which NDCG. Using a hand-written implementation versus
          \texttt{pytrec\_eval} produces a $\sim 2.5$~pt difference on
          TREC-COVID, sufficient on its own to explain a quarter of
          our largest BEIR gap.

    \item \textbf{Per-query min-max normalization before linear
          aggregation}: Eq.~11 of \cite{long2025precise} writes the
          PAGC score as a plain average $\tfrac{1}{|R|+1}\sum$. The
          authors' released code (\texttt{src/utils/aggregate\_utils.py},
          functions \texttt{normalize\_scores} and
          \texttt{aggregate\_linear}) inserts a per-query, per-component
          min-max normalization to $[0,1]$ before the average. Without
          it, RG-YN's $[0,1]$ probabilities and GCCP's contrastive logits
          live on incomparable scales and the aggregation reduces to
          ``whichever component has larger raw magnitude.'' This was
          undocumented in the paper but is reproduced verbatim from the
          released code.
\end{enumerate}

\paragraph{Practical implication.} A reader who tries to reproduce
this method from the paper alone will fail. Each of these knobs
fails \emph{silently}: the model produces seemingly reasonable
probabilities, the pipeline runs to completion, and the final number
is just wrong. We argue (Section~\ref{sec:discussion}) that this is
the broadly relevant lesson, not the specific knobs. The general
phenomenon -- IR papers under-document the engineering layer of the
implementation -- has been raised before in the IR-reproducibility
literature~\cite{breuer2020repro,lin2022leveraging,lin2018msmarco};
our contribution here is a concrete recent instance with a published
catalogue.

\paragraph{Reproducibility hazards we surfaced while pursuing a
spurious gap.} An earlier draft of this paper compared our
DL20/Flan-T5-Large PAGC ($0.6515$) against $0.6910$, which we had
recorded as the paper's reported value, and treated the resulting
$5.7\%$ gap as the largest residual to explain. After re-extracting
the paper's Table~1 RG-YN+GCCP row directly from the PDF, the
correct paper-side value is $0.6281$ -- meaning we actually
\emph{exceed} the paper on this cell by $+3.7\%$. The $0.6910$
appears nowhere in the paper. We present this as a cautionary
finding for reproducibility-track audits: the source of every
``paper'' number a reproduction reports must be traceable to a
specific table cell, and we did not always meet that bar in
intermediate drafts.

The diagnostic steps we ran while pursuing the spurious gap
nonetheless surfaced reproducibility-relevant facts. Switching
from sparse \texttt{scipy.sparse.linalg.eigsh} to dense
\texttt{numpy.linalg.eigh} (the authors' choice) moves PAGC by
$0.1$~pt -- below noise. Porting the authors' exact NLTK-based
sentence segmentation lifts DL20/T5-Large GCCP by $+1.2$~pts and
PAGC by $+0.5$~pts (PAGC then exceeds the paper by $+4.5\%$).
Re-running with five seeds $\{0, 42, 929, 12345, 2023\}$ gives
identical NDCG@10 across seeds (std $=0.0000$), so seed plays no
role in the pipeline output. A future re-implementer with a real
DL20 gap can rule out these three causes without re-running.

% ============================================================
\section{Statistical Significance}
% ============================================================
\label{sec:stats}

The original paper reports per-cell paired $t$-tests at $p \leq 0.05$
for ``\,$+$~GCCP vs single component'' only (the dagger markers in
Tables~1--2 of \cite{long2025precise}); it does not report
multiple-comparison correction, family-wise error control, PAGC vs
GCCP-alone tests, or any of the BEIR-E5 / decoder-only / DBPedia-Entity
analyses we add. We extend the paper's per-cell t-tests with paired
bootstrap two-sided $p$-values (\textbf{10{,}000 resamples}, seed 929)
on per-query NDCG@10 deltas, and 95\% CIs on the mean delta, for
\emph{three} comparison families (PAGC vs RG-YN, PAGC vs GCCP, GCCP
vs RG-YN) across all
\textbf{22 primary} (dataset $\times$ retrieval $\times$ model) settings
(14 TREC-DL cells = 7 (backbone $\times$ retrieval) combinations
$\times$ \{DL19, DL20\}; the 7 combinations are T5-Large/BM25,
T5-XL/BM25, T5-XL/E5, LLaMA-3.1-8B/BM25, Qwen-2.5-7B/BM25,
Mistral-7B/BM25, Qwen-2.5-72B-AWQ/BM25; plus 8 BEIR cells under
Flan-T5-XL/E5).
The full enumeration with $n_q$ and per-comparison numbers is in
Appendix~\ref{app:full-stats}. Because we report 66 pairwise tests
in total, we additionally apply \textbf{Holm-Bonferroni~\cite{holm1979} step-down
correction} within each comparison family ($k{=}22$ tests per family)
to control the family-wise error rate; corrected p-values are denoted
$p_{\text{Holm}}$. We apply Holm \emph{within each comparison family
of $k=22$} rather than across the joint $k=66$ union of all three
families because the three central comparisons (PAGC vs RG-YN, PAGC
vs GCCP, GCCP vs RG-YN) test conceptually distinct hypotheses --
the paper's central claims are family-conditional (e.g.\ ``GCCP
contributes a per-query signal beyond RG-YN'' vs ``aggregating GCCP
with RG-YN beats either alone''), and pooling across families
discards information about which hypothesis is supported. Per-query
scores for all settings are released so any reader can recompute
under the joint $k=66$ correction or any other procedure (e.g.\
Benjamini-Hochberg).
Auxiliary settings (the 5-seed determinism check on DL20/T5-Large/BM25,
and the authors' MDS-port sentence-segmentation cell) re-test the same data and
are excluded from the Holm family to avoid inflating $k$ without adding
independent hypotheses. Table~\ref{tab:stats} shows representative settings.

\begin{table}[t]
\centering
\small
\begin{tabular}{llcccc}
\toprule
\textbf{Setup} & \textbf{Comparison} & \textbf{$\Delta$} & \textbf{$p_{\text{raw}}$} & \textbf{$p_{\text{Holm}}$} & \\
\midrule
\multirow{3}{*}{DL19 / T5-L / BM25}
 & PAGC vs RG-YN & $+0.0219$ & $0.119$ & $0.582$ & ns \\
 & PAGC vs GCCP  & $+0.0512$ & $0.000$ & $0.000$ & *** \\
 & GCCP vs RG-YN & $-0.0293$ & $0.145$ & $1.000$ & ns \\
\midrule
\multirow{3}{*}{DL20 / T5-L / BM25}
 & PAGC vs RG-YN & $+0.0386$ & $0.013$ & $0.132$ & ns \\
 & PAGC vs GCCP  & $+0.0369$ & $0.000$ & $0.000$ & *** \\
 & GCCP vs RG-YN & $+0.0016$ & $0.939$ & $1.000$ & ns \\
\midrule
\multirow{3}{*}{DL19 / T5-XL / E5}
 & PAGC vs RG-YN & $+0.0433$ & $0.004$ & $0.048$ & * \\
 & PAGC vs GCCP  & $+0.0095$ & $0.486$ & $1.000$ & ns \\
 & GCCP vs RG-YN & $+0.0338$ & $0.112$ & $1.000$ & ns \\
\midrule
\multirow{3}{*}{DL19 / Qwen-2.5-7B}
 & PAGC vs RG-YN & $+0.0685$ & $0.000$ & $0.000$ & *** \\
 & PAGC vs GCCP  & $+0.0173$ & $0.073$ & $0.952$ & ns \\
 & GCCP vs RG-YN & $+0.0512$ & $0.010$ & $0.144$ & ns \\
\midrule
\multirow{3}{*}{DL19 / Qwen-72B-AWQ}
 & PAGC vs RG-YN & $+0.0875$ & $0.000$ & $0.000$ & *** \\
 & PAGC vs GCCP  & $+0.0142$ & $0.056$ & $0.784$ & ns \\
 & GCCP vs RG-YN & $+0.0733$ & $0.001$ & $0.012$ & * \\
\midrule
\multirow{3}{*}{BEIR / SciFact / E5}
 & PAGC vs RG-YN & $+0.0597$ & $0.000$ & $0.000$ & *** \\
 & PAGC vs GCCP  & $+0.0150$ & $0.189$ & $1.000$ & ns \\
 & GCCP vs RG-YN & $+0.0446$ & $0.009$ & $0.144$ & ns \\
\midrule
\multirow{3}{*}{BEIR / NFCorpus / E5}
 & PAGC vs RG-YN & $+0.0172$ & $0.000$ & $0.000$ & *** \\
 & PAGC vs GCCP  & $+0.0048$ & $0.229$ & $1.000$ & ns \\
 & GCCP vs RG-YN & $+0.0124$ & $0.041$ & $0.580$ & ns \\
\midrule
\multirow{3}{*}{BEIR / DBPedia / E5}
 & PAGC vs RG-YN & $+0.0808$ & $0.000$ & $0.000$ & *** \\
 & PAGC vs GCCP  & $-0.0144$ & $0.002$ & $0.032$ & * \\
 & GCCP vs RG-YN & $+0.0953$ & $0.000$ & $0.000$ & *** \\
\bottomrule
\end{tabular}
\caption{Paired bootstrap on NDCG@10 (10{,}000 resamples, seed 929)
with Holm-Bonferroni correction across the 22 primary settings within
each comparison family. Significance markers use $p_{\text{Holm}}$.
\textbf{ns}: not significant at $p_{\text{Holm}}<0.05$;
*: $p_{\text{Holm}}<0.05$;
**: $p_{\text{Holm}}<0.01$;
***: $p_{\text{Holm}}<0.001$. Note the BEIR/DBPedia row: PAGC is
Holm-significantly \emph{worse} than GCCP-alone
($\Delta{=}{-}0.0144$, $p_{\text{Holm}}{=}0.032$) -- aggregating
RG-YN with GCCP under E5 retrieval can hurt at large $n_q$.}
\label{tab:stats}
\end{table}

\paragraph{Findings (across the 22 primary settings).}
(i) \textbf{PAGC vs RG-YN is Holm-significant in 12/22 settings,
all directionally positive} -- 6/8 BEIR-E5 sets clear the bar
(only Touch\'e-2020 $n_q{=}49$ and TREC-News $\Delta{=}{+}0.007$
do not), plus 6/14 TREC-DL cells skewed toward larger backbones
and the single E5 cell.
(ii) \textbf{PAGC vs GCCP-alone is Holm-significant in 5/22 with
mixed signs:} 4 positive (all BM25 setups) and \textbf{1 negative}
(BEIR/DBPedia-Entity under E5, $\Delta{=}{-}0.0144$,
$p_{\text{Holm}}{=}0.032$). All E5 cells are either ns or
negative -- aggregating RG-YN with GCCP under a strong dense
retriever stops being uniformly beneficial.
(iii) \textbf{GCCP-alone vs RG-YN is Holm-significant in only
3/22:} the two Qwen-72B-AWQ cells and BEIR/DBPedia-Entity --
either the largest decoder backbone or the largest BEIR set. In
the other 19 settings the contrastive-anchor mechanism does not
produce a detectable per-query signal, even though aggregate
means improve.

\paragraph{Power caveat.} On TREC-DL with $n{=}43$/$54$, ``not
significant after Holm'' is consistent with both ``no effect'' and
``a real effect we cannot detect.'' We therefore make claims of
\emph{absence of evidence}, not \emph{evidence of absence}, for
small-$n$ TREC-DL cells; per-query scores are released for
re-analysis at higher $n$. The implication remains more nuanced
than ``GCCP+aggregation always wins'': under E5, aggregation is
never Holm-significantly better than GCCP-alone, and on the largest
BEIR-E5 set it is significantly \emph{worse}.

% ============================================================
\section{Beyond the Paper}
\label{sec:beyond}
% ============================================================

\subsection{E5 dense retrieval}
\label{sec:e5}
The original paper uses BM25 as the only first-stage retriever. We
replaced it with E5 (\texttt{intfloat/e5-base-v2})~\cite{wang2022e5}
keeping the rest of the pipeline unchanged.

\begin{table}[t]
\centering
\small
\begin{tabular}{lcccccc}
\toprule
 & \multicolumn{2}{c}{\textbf{First-stage}}
 & \multicolumn{2}{c}{\textbf{GCCP}}
 & \multicolumn{2}{c}{\textbf{PAGC}} \\
\cmidrule(lr){2-3}\cmidrule(lr){4-5}\cmidrule(lr){6-7}
\textbf{Dataset} & BM25 & E5 & +BM25 & +E5 & +BM25 & +E5 \\
\midrule
\multicolumn{7}{l}{\emph{TREC DL}} \\
DL19         & 0.5058 & 0.7086 & 0.6844 & \textbf{0.7090} & 0.7030 & \textbf{0.7185} \\
DL20         & 0.4796 & 0.7051 & 0.6668 & \textbf{0.7069} & 0.6760 & \textbf{0.7177} \\
\multicolumn{7}{l}{\emph{BEIR}} \\
SciFact      & 0.6789 & 0.7274 & 0.6815 & \textbf{0.7026} & 0.6840 & \textbf{0.7176} \\
NFCorpus     & 0.3218 & 0.3517 & 0.3758 & \textbf{0.3836} & 0.3728 & \textbf{0.3884} \\
\bottomrule
\end{tabular}
\caption{NDCG@10 with Flan-T5-XL, BM25 vs.\ E5-base-v2 first-stage.}
\label{tab:e5-dl}
\end{table}

E5 lifts the first-stage by $+3$ to $+23$~pts over BM25 across these
four datasets (DL19/DL20 see large $+20$/$+23$~pt lifts; the BEIR
sets see smaller $+3$ to $+5$~pt lifts). On TREC DL, our PAGC$+$E5
exceeds \emph{both} the paper's 2-component RG-YN+GCCP (Table~1,
T5-XL: $0.6969$ on DL19, $0.6810$ on DL20) by $+2.2$/$+3.7$~pts and
the paper's 3-component PAGC-QSG (Table~2, T5-XL: $0.7068$ on DL19,
$0.6872$ on DL20) by $+1.2$/$+3.0$~pts respectively. On SciFact,
PAGC$+$E5 ($0.7176$) is $+3.4$~pts over our PAGC$+$BM25 ($0.6840$)
and $+3.2$~pts over the paper's T5-XL Table~1 RG-YN+GCCP number
($0.6854$).

\paragraph{Negative result: PAGC can hurt over a strong first-stage.}
The marginal contribution of GCCP re-ranking on top of E5 is much
smaller than on top of BM25, and in one cell it is negative:
\begin{itemize}[leftmargin=1.5em,nosep]
    \item DL20: PAGC$+$BM25 buys $+0.197$ over BM25; PAGC$+$E5 buys
          $+0.013$ over E5.
    \item DL19: PAGC$+$BM25 buys $+0.197$ over BM25; PAGC$+$E5 buys
          $+0.010$ over E5.
    \item SciFact: PAGC$+$BM25 buys $+0.005$ (essentially nothing);
          \textbf{PAGC$+$E5 buys $-0.010$, i.e.\ PAGC slightly
          \emph{hurts} over E5 alone}.
    \item NFCorpus: PAGC$+$BM25 buys $+0.051$; PAGC$+$E5 buys
          $+0.037$ -- here the dense first-stage is itself weaker
          ($0.3517$ vs BM25's $0.3218$) and the reranker continues
          to add signal.
\end{itemize}
\textbf{The contrastive anchor is most useful when the candidate
list is noisy.} When the first-stage already returns a clean list
(SciFact + E5), GCCP rearrangement has no additional signal to
extract and can introduce noise. This is consistent with our
statistical-significance findings (Section~\ref{sec:stats}): on E5
setups, PAGC vs GCCP is never Holm-significant; on BM25 setups it
is significant in 4/9.

\subsection{Decoder-only LLMs}
\label{sec:decoder}
The original paper covers only Flan-T5 / Flan-UL2, all encoder-decoder.
We adapt the RG-YN and GCCP prompts to chat-template format and run
four decoder-only models: LLaMA-3.1-8B-Instruct~\cite{llama3report},
Qwen-2.5-7B-Instruct~\cite{qwen25technical},
Mistral-7B-Instruct-v0.3~\cite{jiang2023mistral}, and
Qwen-2.5-72B-Instruct (AWQ-quantized) for a scaling data point. We
prompt via
\texttt{tokenizer.apply\_chat\_template(\dots, add\_generation\_prompt=True)},
reading next-token probabilities at the assistant turn boundary and
aggregating over case/space variants of the target tokens.

\paragraph{Implementation note.} The GCCP prompt asks the model to
``Output Passage A or Passage B''; an instruction-tuned chat model
emits ``Passage'' first, not ``A''. We therefore prime the assistant
turn with the literal string \texttt{`Passage '} before reading next-token
logits, so we score $P(\text{`A'} | \dots\text{Passage }) /
(P(\text{`A'}) + P(\text{`B'}))$ -- exactly mirroring the T5
\texttt{decoder\_input\_text='<pad> Passage '} convention from the
authors' code. Without this primer, the GCCP score collapses (we
verified a 4× reduction in NDCG@10 on a Qwen-2.5-0.5B smoke test).

\begin{table}[t]
\centering
\small
\begin{tabular}{llcccc}
\toprule
\textbf{Dataset} & \textbf{Model} & \textbf{BM25} & \textbf{RG-YN} & \textbf{GCCP} & \textbf{PAGC} \\
\midrule
DL19 & LLaMA-3.1-8B-Instruct    & 0.5058 & 0.6427 & 0.6521 & 0.6723 \\
DL19 & Qwen-2.5-7B-Instruct     & 0.5058 & 0.6527 & 0.7039 & 0.7212 \\
DL19 & Mistral-7B-Instr.\ v0.3  & 0.5058 & 0.5716 & 0.5617 & 0.6429 \\
DL19 & Qwen-2.5-72B-AWQ         & 0.5058 & 0.6590 & \textbf{0.7323} & \textbf{0.7465} \\
\midrule
DL20 & LLaMA-3.1-8B-Instruct    & 0.4796 & 0.5795 & 0.5866 & 0.6166 \\
DL20 & Qwen-2.5-7B-Instruct     & 0.4796 & \textbf{0.6360} & 0.6447 & 0.6641 \\
DL20 & Mistral-7B-Instr.\ v0.3  & 0.4796 & 0.5360 & 0.5105 & 0.5983 \\
DL20 & Qwen-2.5-72B-AWQ         & 0.4796 & 0.6311 & \textbf{0.6983} & \textbf{0.6980} \\
\bottomrule
\end{tabular}
\caption{NDCG@10 on TREC DL with decoder-only LLMs adapted via chat
templates (novel; not in the original paper). For comparison: Flan-T5-Large
(780M) PAGC = 0.6834 / 0.6515 on DL19/DL20; Flan-T5-XL (3B) = 0.7030 /
0.6760; Flan-UL2 (20B) = 0.7095 / 0.7009. The original paper reports
Flan-UL2 RG-YN+GCCP = 0.7206 / 0.7023 (Table~1 of \cite{long2025precise}).
\textbf{Qwen-2.5-72B-AWQ on DL19 (0.7465) is the strongest cell in
this table, surpassing both our reproduced Flan-UL2 and the paper's
own reported Flan-UL2 number on DL19} -- a $\sim$36~GB AWQ-quantized
open-weight model running on a single 48~GB GPU does so without
matching the paper's compute budget.}
\label{tab:decoder}
\end{table}

\paragraph{Findings.} The four decoder-only models span the full
range from \emph{below} Flan-T5-Large (780M) to \emph{above} the
paper's reported Flan-UL2 (20B):

\begin{itemize}[leftmargin=1.5em,nosep]
    \item \textbf{Qwen-2.5-72B-AWQ is the strongest model we
          tested.} On DL19 PAGC it reaches $0.7465$, surpassing
          both our reproduced Flan-UL2 ($0.7095$) and the paper's
          Flan-UL2 ($0.7206$, RG-YN+GCCP) by $+2.59$~pts. On DL20
          PAGC it sits slightly below the paper's Flan-UL2
          ($0.6980$ vs $0.7023$, $-0.43$~pts) but matches our
          reproduced Flan-UL2 within $0.003$ -- a 4-bit AWQ
          open-weight 72B on a single 48~GB GPU competitive with
          the paper's best dense-FP16 result.
    \item \textbf{Cross-architecture transfer onto BEIR.} On a
          compute-capped 50-query subset ($\sim$4~min/query at 72B):
          SciFact reproduces the strict ordering
          $\text{E5}(0.7993) < \text{RG-YN}(0.8074) <
          \text{GCCP}(0.8586) < \text{PAGC}(\mathbf{0.8734})$;
          \textbf{DBPedia-Entity reproduces the negative finding from
          Section~\ref{sec:beir-e5}}:
          $\text{RG-YN}(0.3667) < \text{PAGC}(0.4139) <
          \text{E5}(0.4168) < \text{GCCP}(\mathbf{0.4380})$ -- PAGC
          underperforms GCCP-alone \emph{and} the E5 first-stage,
          surviving a complete change of backbone family
          (encoder-decoder $\to$ decoder-only), parameter count
          ($3$B $\to$ $72$B), and quantization regime
          (FP16 $\to$ 4-bit AWQ).
    \item At 7--8B, \textbf{Qwen-2.5-7B} is competitive with
          Flan-T5-XL (DL19 PAGC $0.7212$ vs $0.6939$, paired
          bootstrap $\Delta{=}{+}0.027$, $p{=}0.125$, not significant
          on $n{=}43$); \textbf{LLaMA-3.1-8B} is below Flan-T5-Large
          ($-1.1$/$-3.5$~pts on DL19/DL20) despite being $10\times$
          larger by parameter count; \textbf{Mistral-7B-v0.3}
          underperforms (especially on GCCP, A/B discrimination
          $0.5617$/$0.5105$).
\end{itemize}

\noindent Two complementary conclusions: at fixed 7--8B scale,
backbone family matters at least as much as parameter count (the
Qwen-2.5 vs Mistral gap $+0.078$ on DL19 PAGC exceeds Flan-T5-Large
vs Flan-UL2 at $+0.026$); and the contrastive anchor transfers
cleanly to large quantized decoder-only models.

\subsection{Aggregation methods}
\label{sec:agg}
We compare the paper's equal-weight linear aggregation against an
$\alpha$-weighted linear sweep, and Borda / Condorcet / Copeland
rank voting, across 8 (dataset $\times$ retrieval $\times$ model)
cells. Because we use the 2-component PAGC = RG-YN $+$ GCCP
throughout, $|R|=2$, and Borda / Condorcet / Copeland are
mathematically identical in this regime; we therefore cannot
differentiate among positional aggregation methods within our
setup. The base paper's Section~4.4.1 study uses the 3-component
PAGC-QYG and PAGC-QSG ($|R|=3$), where these methods can in
principle diverge; the paper's finding that they nonetheless
yield ``nearly identical NDCG@10'' is a substantive observation
about that 3-component regime, not a methodological artifact.
Table~\ref{tab:agg} reports our full sweep.

\begin{table*}[t]
\centering
\small
\begin{tabular}{lcccccc}
\toprule
\textbf{Setting} & RG-YN only & GCCP only & lin $\alpha{=}0.5$ (paper) & lin $\alpha{=}0.25$ & lin $\alpha{=}0.75$ & Borda \\
\midrule
DL19 / T5-Large / BM25       & 0.6634 & 0.6341 & 0.6852 & 0.6730 & 0.6848 & \textbf{0.6966} \\
DL19 / T5-XL    / BM25       & 0.6737 & 0.6747 & \textbf{0.6939} & 0.6881 & 0.6933 & 0.6936 \\
DL19 / T5-XL    / E5         & 0.6752 & 0.7090 & 0.7185 & \textbf{0.7267} & 0.7026 & 0.7154 \\
DL20 / T5-Large / BM25       & 0.6121 & 0.6137 & 0.6507 & 0.6399 & 0.6488 & \textbf{0.6510} \\
DL20 / T5-XL    / BM25       & 0.6512 & 0.6701 & 0.6725 & 0.6768 & 0.6675 & \textbf{0.6865} \\
DL20 / T5-XL    / E5         & 0.6789 & 0.7069 & 0.7177 & \textbf{0.7199} & 0.7167 & 0.7113 \\
BEIR-SciFact  / T5-XL / E5   & 0.6579 & 0.7026 & 0.7176 & \textbf{0.7238} & 0.6862 & 0.6915 \\
BEIR-NFCorpus / T5-XL / E5   & 0.3712 & 0.3836 & 0.3884 & \textbf{0.3917} & 0.3817 & 0.3892 \\
\midrule
\textbf{Avg.\ (over 8 cells)} & 0.6230 & 0.6368 & 0.6556 & 0.6550 & 0.6477 & 0.6544 \\
\bottomrule
\end{tabular}
\caption{Aggregation ablation across 8 (dataset $\times$ retrieval $\times$ model)
cells. Bold marks the best aggregation per cell among the four
non-baseline rows. Borda, Condorcet, and Copeland are numerically
identical at $|R|=2$.}
\label{tab:agg}
\end{table*}

\paragraph{Findings (across the 8-cell sweep).}
(i) Borda is not consistently better than the paper's
$\alpha{=}0.5$ linear: it wins in 3/8 cells; the cell-mean over
all 8 is essentially tied ($0.6544$ vs $0.6556$). The $+1.1$~pt
gap on DL19/T5-Large/BM25 we initially highlighted does not
generalize.
(ii) $\alpha{=}0.25$ wins on every E5 cell (4/4) and on
DL20/T5-XL/BM25, but loses on BM25 with smaller models; cell-mean
$0.6550$ vs $0.6556$. The optimal-$\alpha$ \emph{direction} is
thus retrieval-quality-dependent (E5 wants more GCCP weight) but
the magnitudes are small enough that $\alpha{=}0.5$ is a robust
default.
(iii) At $|R|=2$, Borda / Condorcet / Copeland are mathematically
identical, so we cannot differentiate them; the paper's
Section~4.4.1 study at $|R|=3$ remains the relevant evidence.

\subsection{Beyond the two-component PAGC: adding RG-S}
\label{sec:pagc_qsg}

The paper's strongest reported variant is
PAGC-QSG~$=$~QG~$+$~RG-S~$+$~GCCP. We add the RG-S$(0,4)$ component
and report the 3-component PAGC-RS-YN-GCCP variant; QG is left for
follow-up since it requires full-sequence query log-likelihood.
RG-S$(0,4)$ follows the authors' \texttt{clf\_ranking.py} template~4
(prompt ``From a scale of 0 to 4, judge the relevance''; score is
$P(\text{`4'})$ via the softmax of the 5 integer labels;
truncation 128 tokens; T5 decoder primer \texttt{`<pad> '}).

\begin{table}[t]
\centering
\small
\begin{tabular}{lcc|cc}
\toprule
& \multicolumn{2}{c}{\textbf{Ours (T5-Large/BM25)}}
& \multicolumn{2}{c}{\textbf{Paper (Table 4 hom.)}} \\
\textbf{Method}        & DL19 & DL20 & TREC avg & BEIR avg \\
\midrule
RG-S only              & 0.5872 & 0.5204 & --     & --     \\
RG-YN only             & 0.6634 & 0.6121 & --     & --     \\
GCCP only              & 0.6341 & 0.6137 & --     & --     \\
PAGC = RG-YN $+$ GCCP  & 0.6852 & 0.6507 & --     & --     \\
PAGC = RG-S $+$ GCCP   & 0.6547 & 0.6058 & --     & --     \\
PAGC-RS-YN-GCCP        & 0.6856 & 0.6325 & 0.6634 & 0.4538 \\
\bottomrule
\end{tabular}
\caption{PAGC-RS-YN-GCCP (paper's Table~4 ``RG-S+RG-YN+GCCP''
homogeneous variant) on Flan-T5-Large/BM25. Our TREC average
$(0.6856 + 0.6325)/2 = 0.6591$ vs the paper's $0.6634$ -- a
$0.4\%$ gap, consistent with our other reproductions.}
\label{tab:pagc_rs}
\end{table}

\paragraph{Findings.} (a) The 3-component PAGC reproduces the paper's
Table~4 number within $0.4\%$ on TREC average. (b) On DL19 the
three-component aggregation is essentially tied with the two-component
$(0.6856$ vs $0.6852)$. (c) \textbf{On DL20 the three-component
aggregation \emph{loses} to the two-component} ($0.6325$ vs $0.6507$,
$-1.8$ pts). RG-S alone is much weaker than RG-YN on DL20 ($0.5204$
vs $0.6121$); aggregating a weak component drags the mean rank.
This is a real but subtle phenomenon hidden by the paper's choice to
report only TREC$+$BEIR averaged; per-dataset, the value of adding
RG-S is non-monotonic. The paper's strongest variant (PAGC-QSG)
remains untested by us; we expect QG to behave similarly to the
other components -- helpful when its component scores are uncorrelated
with RG-YN/GCCP errors, but not always so.

\subsection{Position relative to listwise/pairwise baselines}
\label{sec:rankgpt}

To contextualize our 2-component PAGC against the listwise and
pairwise baselines the paper reports in its Table~2, we summarize
the paper-reported numbers (no re-run) alongside our PAGC numbers
at the same Flan-T5 scales.

\begin{table}[t]
\centering
\small
\begin{tabular}{lcc}
\toprule
\textbf{Method (Flan-T5-Large, DL19/DL20)} & DL19 & DL20 \\
\midrule
BM25 (no rerank)                       & 0.5058 & 0.4796 \\
\midrule
\multicolumn{3}{l}{\emph{Paper-reported (Table~2 of \cite{long2025precise}):}}\\
\quad RankGPT                          & 0.6690 & 0.6260 \\
\quad PRP-Allpair                      & 0.6646 & 0.6201 \\
\quad PRP-Heapsort                     & 0.6570 & 0.6190 \\
\quad PRP-Graph-40                     & 0.6690 & 0.6270 \\
\quad Setwise-Heapsort                 & 0.6700 & 0.6180 \\
\quad PAGC-QSG (paper, 3-comp)         & 0.6973 & 0.6433 \\
\midrule
\multicolumn{3}{l}{\emph{Ours (PAGC = RG-YN $+$ GCCP, 2-comp):}}\\
\quad PAGC / Flan-T5-Large             & \textbf{0.6852} & \textbf{0.6507} \\
\quad PAGC / Flan-T5-XL                & \textbf{0.6939} & 0.6760 \\
\bottomrule
\end{tabular}
\caption{Position of our 2-component PAGC relative to the listwise
and pairwise baselines the paper reports in its Table~2 (no
re-run; numbers cited verbatim). Our PAGC sits inside the bracket
of paper-reported listwise/pairwise baselines -- slightly above
all of the paper's listwise/pairwise rows and slightly below its
3-component PAGC-QSG.}
\label{tab:rankgpt}
\end{table}

\paragraph{Position.} Our 2-component PAGC ($0.6852$ on DL19/T5-Large)
sits inside the bracket of the paper's reported listwise/pairwise
methods -- slightly above all five (RankGPT $0.6690$,
PRP-Allpair $0.6646$, PRP-Heapsort $0.6570$, PRP-Graph-40 $0.6690$,
Setwise-Heapsort $0.6700$) and slightly below the paper's
3-component PAGC-QSG ($0.6973$). At Flan-T5-XL the picture is
similar but the listwise/pairwise baselines are stronger:
PRP-Allpair ($0.7108$) and PRP-Heapsort ($0.7050$) edge above our
2-component PAGC ($0.6939$), and PAGC-QSG ($0.7068$) leads. The
qualitative ordering -- 3-component aggregation $>$ 2-component
aggregation $>$ listwise/pairwise on T5-Large; pairwise variants
become competitive with 2-component aggregation at T5-XL -- is
consistent with the paper's Table~2 narrative.\footnote{We
attempted an \texttt{llm-rankers} re-run of RankGPT-listwise on
the same Flan-T5 backbones to enable a re-implemented head-to-head
comparison; our re-run produces NDCG@10 $\approx 0.5$, well below
the paper's reported RankGPT $\approx 0.6690$. We were unable to
identify the cause from the released artifacts (sliding-window
size, top-$k$ truncation, prompt template, exact RankGPT variant
of~\cite{sun2023chatgpt}) and therefore do \emph{not} include the
re-run as a comparison baseline; the released code is in our
artifact for any reader who wants to take the discrepancy further.}

\subsection{Anchor and parameter ablations}
\label{sec:anchor-ablation}
We ablate the anchor builder against three baselines: random
candidate, top-1 BM25, and a top-3 sentence-interleaved composite.
Full results on the complete query sets (43 queries on DL19,
54 queries on DL20) with Flan-T5-Large appear in
Table~\ref{tab:anchor-full}.

\begin{table}[t]
\centering
\small
\begin{tabular}{lcc|cc}
\toprule
& \multicolumn{2}{c}{\textbf{DL19}} & \multicolumn{2}{c}{\textbf{DL20}} \\
\textbf{Anchor builder} & GCCP & PAGC & GCCP & PAGC \\
\midrule
Random passage   & 0.6394 & 0.6945 & 0.6103 & 0.6474 \\
Top-1 BM25       & \textbf{0.6511} & \textbf{0.6948} & 0.6131 & 0.6485 \\
Top-3 composite  & 0.6410 & 0.6947 & \textbf{0.6280} & \textbf{0.6572} \\
Spectral MDS (paper) & 0.6341 & 0.6852 & 0.6137 & 0.6507 \\
\bottomrule
\end{tabular}
\caption{Anchor-builder ablation on full DL19/DL20 with
Flan-T5-Large. RG-YN is identical across rows by construction
(0.6634 / 0.6121) and omitted. \textbf{Spectral MDS is not the
strongest anchor in our reproduction.}}
\label{tab:anchor-full}
\end{table}

\paragraph{Finding (replicates paper's TREC-DL exception, magnifies it).}
On DL19, Top-1 BM25 (a one-line anchor builder) beats spectral MDS
by $+1.7$~pts on GCCP and $+1.0$~pt on PAGC. On DL20, a Top-3
sentence-interleaved composite wins by $+1.4$~pts and $+0.7$~pts. On
the (DL19$+$DL20)/2 average that the paper uses, our GCCP$+$Top-1 is
$0.6321$ vs GCCP$+$Spectral $0.6239$, a $+0.008$ gap.

The paper itself reports the same direction on TREC DL (Table 5 of
\cite{long2025precise}): GCCP$+$Top $= 0.6099$, GCCP$+$Spectral
$= 0.6076$, a $+0.002$ gap, and explicitly notes: \emph{``The Top
method, while also performing well, generally falls slightly short
of the Spectral method, except in the GCCP model on the TREC DL
benchmark, where it marginally outperforms the Spectral method.''}
Our reproduction therefore \emph{confirms} the paper's reported
TREC-DL exception, but we find the magnitude $\sim 4\times$ larger
(+0.008 vs the paper's +0.002). On the (DL19$+$DL20)/2 PAGC average
the paper's authors report Spectral$>$Top across all PAGC variants
in their Table 5; we have not yet replicated the BEIR portion of
their anchor ablation, which is where they observe the larger
Spectral$>$Top gap.

\paragraph{Implication.} The result is not that spectral MDS does
not work; it is that on TREC DL with GCCP-alone, the cheapest
possible anchor (the top-1 BM25 candidate, no clustering, no MDS) is
already as good. The paper's \emph{aggregate} case for spectral MDS
rests on its BEIR performance (Table~5: $0.4471$ vs Top's $0.4346$
averaged across 8 BEIR sets).

\paragraph{Per-dataset BEIR anchor (full 8 datasets).} We ran the
same 4-anchor ablation on all 8 BEIR sets used in the original
paper. Table~\ref{tab:beir-anchor} reports GCCP NDCG@10 per set;
PAGC follows the same pattern (omitted for space).

\begin{table*}[t]
\centering
\small
\setlength{\tabcolsep}{4pt}
\begin{tabular}{lcccccccc|cc}
\toprule
\textbf{Anchor} & SciFact & NFCorpus & TREC-COVID & Touch\'e & Robust04 & TREC-News & Signal1M & DBPedia & \textbf{Ours avg.} & \textbf{Paper avg.} \\
\midrule
Random           & 0.5781 & 0.3213 & \textbf{0.7701} & \textbf{0.2954} & 0.4386 & \textbf{0.4384} & 0.2894 & 0.4075 & 0.4423 & 0.4253 \\
Top-1 BM25       & 0.6659 & 0.3303 & 0.7609 & 0.2894 & 0.4423 & 0.4218 & 0.2734 & 0.4255 & 0.4512 & 0.4346 \\
Top-3 composite  & \textbf{0.6692} & \textbf{0.3381} & 0.7649 & 0.2869 & \textbf{0.4505} & 0.4261 & \textbf{0.3020} & \textbf{0.4268} & \textbf{0.4581} & --     \\
Spectral (paper) & 0.6195 & 0.3311 & 0.7564 & 0.2694 & 0.4380 & 0.4245 & 0.2990 & 0.4186 & 0.4446 & \textbf{0.4471} \\
\bottomrule
\end{tabular}
\caption{GCCP NDCG@10 across all 8 BEIR datasets with Flan-T5-Large /
BM25, varying only the anchor builder. ``Ours avg.'' is the simple
mean across 8 datasets; ``Paper avg.'' is the corresponding number
from \cite{long2025precise} Table~5 (paper does not include a
``Top-3 composite'' baseline). \textbf{Spectral never wins on any
of the 8 sets and finishes last on 3/8 (TREC-COVID, Touch\'e-2020,
Robust04); on the other 5 it sits in the middle of the four anchors.}
Top-3 composite wins 5/8 (SciFact, NFCorpus, Robust04, Signal1M,
DBPedia); Random wins 3/8 (TREC-COVID, TREC-News, Touch\'e-2020).}
\label{tab:beir-anchor}
\end{table*}

\paragraph{Finding.} On TREC-DL we replicated the paper's reported
exception (GCCP+Top slightly beats GCCP+Spectral); within the
paper's 4-anchor comparison (Random / Top-1 / Top-3 / Spectral)
\emph{on every single one of the 8 BEIR sets we ran},
Spectral MDS loses to at least one simpler anchor, and
\emph{finishes last on 3/8 sets} (TREC-COVID, Touch\'e-2020, and
Robust04 -- where the cheapest possible anchor, Random, beats it).
On the other 5 sets it lands second or third. The aggregate
ordering on our 8 sets (Top-3 $>$ Top-1 $>$ Spectral $>$ Random)
is the \emph{opposite} of the paper's reported aggregate (Spectral
$>$ Top $>$ Random).

\paragraph{Reconciliation: title-only Top closes the gap.}
Our reproduction matches the paper closely on Spectral
($0.4446$ vs $0.4471$, $+0.0025$) but diverges on Top ($+0.017$)
and Random ($+0.017$) in the initial 4-anchor sweep. We tested
four plausible operationalizations of the paper's ``Top'' and
``Random'' (Table~\ref{tab:beir-anchor-variants}). Random-5
averaging closes $\sim 35\%$ of the Random gap (per-seed std
$\leq 0.009$); Top-5 and 64-char Top-1 truncation move further
from the paper; \textbf{Top-1 title-only matches the paper's
``Top'' within $0.003$} ($0.4313$ vs $0.4346$), identifying the
likely operationalization. The residual $\sim 0.011$ on Random
likely reflects a different random-document choice (e.g.\ random
sentence) we have not yet pinned down (see the supplementary
email to the original authors).

\paragraph{Implication for the spectral MDS conclusion.} Even
after matching the paper's likely ``Top'' baseline, the negative
finding on Spectral \emph{strengthens}: across the 9 anchor
variants we tested per dataset, Spectral never finishes first on
any of 8 BEIR sets. It is beaten by Top-3 composite ($0.4581$) and
full-document Top-1 ($0.4512$) on aggregate, so the engineering
cost of spectral MDS is not justified relative to a top-3
sentence-interleaved composite even under the paper's own ``Top''.

\begin{table*}[t]
\centering
\small
\setlength{\tabcolsep}{4pt}
\begin{tabular}{lrrrrr|l}
\toprule
\textbf{Variant} & \textbf{Avg.} & SciFact & DBPedia & Touch\'e & Signal1M
& \textbf{vs paper Table~5} \\
\midrule
Random-1 (our default)        & 0.4423 & 0.5781 & 0.4075 & 0.2954 & 0.2894 & paper Random $0.4253$, $\Delta{+}0.017$ \\
Random-5 averaging            & 0.4364 & 0.5673 & 0.4115 & 0.2805 & 0.2825 & $\Delta{+}0.011$ (closes $\sim 35\%$) \\
\midrule
Top-1 BM25 (full doc, default) & 0.4512 & 0.6659 & 0.4255 & 0.2894 & 0.2734 & paper Top $0.4346$, $\Delta{+}0.017$ \\
Top-3 sentence composite      & \textbf{0.4581} & 0.6692 & 0.4268 & 0.2869 & 0.3020 & not in paper \\
Top-5 sentence composite      & 0.4547 & 0.6476 & 0.4259 & 0.2888 & 0.3029 & not in paper \\
Top-1 64-char truncation      & 0.4157 & 0.5798 & 0.3815 & 0.2653 & 0.2671 & undershoots by $\Delta{-}0.019$ \\
\textbf{Top-1 title-only}     & \textbf{0.4313} & 0.6499 & 0.3776 & 0.2728 & 0.2750 & \textbf{$\Delta{-}0.003$ (matches)} \\
\midrule
Spectral MDS (paper builder)  & 0.4446 & 0.6195 & 0.4186 & 0.2694 & 0.2990 & paper $0.4471$, $\Delta{-}0.003$ \\
\bottomrule
\end{tabular}
\caption{Anchor-operationalization sweep on all 8 BEIR sets,
Flan-T5-Large/BM25 (4 representative datasets shown for the
per-cell columns; full 8-cell results in the released artifact).
\textbf{Top-1 title-only matches the paper's reported ``Top''
average ($0.4313$ vs $0.4346$, $\Delta{-}0.003$)}, identifying the
likely operationalization the paper uses. Random-5 averaging
partially explains the Random gap. Spectral matches the paper
within $0.003$ regardless of the other variants tested.}
\label{tab:beir-anchor-variants}
\end{table*}

\paragraph{Hyperparameter sweep ($m, z, \theta$).} We additionally
sweep $m \in \{5, 10, 15, 20\}$, $z \in \{5, 10, 15, 20\}$, and
$\theta \in \{0.1, 0.2, 0.3, 0.4\}$ on the same DL20/T5-Large/BM25
cell. PAGC NDCG@10 varies by $\pm 1.5$~pts across $m$ and $\theta$
-- a non-trivial range when state-of-the-art deltas on these
benchmarks are themselves $\sim 2$~pts, warranting picking $m$ and
$\theta$ on a held-out split rather than copying the paper's
defaults verbatim. The $z \geq 10$ plateau is explained by the
encoder's 512-token limit: beyond $z=10$, additional anchor
sentences are truncated. The paper's choice $z=10$ is therefore
not fragile, but is also not load-bearing.

\subsection{8-set BEIR generalization under E5 retrieval}
\label{sec:beir-e5}

To test whether the PAGC story holds under a stronger first-stage
retriever, we re-ran the rerank stack with E5-base-v2 dense retrieval
(top-100, FAISS~\cite{johnson2019faiss} \texttt{IndexFlatIP}) on all 8 BEIR datasets used in
the original paper, and reranked with Flan-T5-XL.
Table~\ref{tab:beir-e5} reports NDCG@10 for RG-YN, GCCP, and PAGC,
along with the paired-bootstrap delta (vs RG-YN) and the
Holm-corrected significance marker (10{,}000 resamples, seed~929).
The 8 BEIR-E5 cells are part of the same 22-setting Holm family
described in Section~\ref{sec:stats} -- there is no separate BEIR-E5
correction; the markers shown match Appendix~\ref{app:full-stats}
exactly.

\begin{table*}[t]
\centering
\small
\setlength{\tabcolsep}{4pt}
\begin{tabular}{lcccccc}
\toprule
\textbf{Dataset} & \textbf{$n_q$} & \textbf{RG-YN} & \textbf{GCCP} & \textbf{PAGC}
& \textbf{$\Delta$ PAGC$-$GCCP} & \textbf{$\Delta$ PAGC$-$RG-YN} \\
\midrule
SciFact          & 300 & 0.6579 & 0.7026 & \textbf{0.7176} & $+0.0150$~ns   & $+0.0597$~*** \\
NFCorpus         & 323 & 0.3712 & 0.3836 & \textbf{0.3884} & $+0.0048$~ns   & $+0.0172$~*** \\
TREC-COVID       &  50 & 0.7117 & 0.7752 & \textbf{0.7859} & $+0.0107$~ns   & $+0.0742$~**  \\
Touch\'e-2020    &  49 & 0.2118 & 0.2538 & \textbf{0.2594} & $+0.0056$~ns   & $+0.0477$~ns  \\
DBPedia-Entity   & 400 & 0.3707 & \textbf{0.4659} & 0.4515 & $-0.0144$~\textbf{*} & $+0.0808$~*** \\
TREC-News        &  57 & 0.4346 & 0.4391 & \textbf{0.4415} & $+0.0024$~ns   & $+0.0069$~ns  \\
Robust04         & 249 & 0.4920 & 0.5184 & \textbf{0.5255} & $+0.0071$~ns   & $+0.0335$~*** \\
Signal-1M        &  97 & 0.2334 & \textbf{0.2638} & 0.2612 & $-0.0025$~ns   & $+0.0279$~**  \\
\midrule
\textbf{Avg.\ (8 sets)} & -- & 0.4354 & 0.4753 & \textbf{0.4789} & $+0.0036$ & $+0.0435$ \\
\bottomrule
\end{tabular}
\caption{8-set BEIR generalization under E5 retrieval +
Flan-T5-XL rerank. Holm-correction p-values match
Appendix~\ref{app:full-stats}: the 8 BEIR-E5 cells are part of the
same 22-setting primary Holm family used in
Section~\ref{sec:stats} ($k=22$ tests per comparison family).
\textbf{ns}: $p_{\text{Holm}}\geq 0.05$; *: $<0.05$;
**: $<0.01$; ***: $<0.001$. PAGC beats RG-YN on every set
directionally and Holm-significantly in 6/8 sets. PAGC beats
GCCP-alone in 6/8 sets directionally; the only Holm-significant
difference is on DBPedia-Entity, where \emph{PAGC is significantly
worse than GCCP-alone} ($\Delta=-0.0144$, $p_{\text{Holm}}{=}0.032$).}
\label{tab:beir-e5}
\end{table*}

\paragraph{Findings.}
(i) \textbf{PAGC vs RG-YN holds up under E5}: Holm-sig in 6/8,
directional in 8/8 (the 2 non-sig sets are the smallest by query
count, a power problem).
(ii) \textbf{PAGC vs GCCP-alone is the surprising negative}: in
\emph{none} of 8 sets is PAGC Holm-significantly better than
GCCP-alone, and on the largest (DBPedia-Entity, $n{=}400$) PAGC is
Holm-significantly \emph{worse}. Aggregating with RG-YN does not
consistently help once the first-stage retriever is strong --
echoing our DL/E5 finding (Table~\ref{tab:agg}) that the optimal
$\alpha$ shifts toward GCCP under E5.
(iii) \textbf{GCCP-alone vs RG-YN}: Holm-sig in 1/8
(DBPedia-Entity ***), directional in 8/8 -- the same picture as
TREC-DL but with DBPedia standing out as the E5 cell where the
contrastive-anchor mechanism is decisively better than yes/no
scoring.

% ============================================================
\section{Efficiency}
% ============================================================
\label{sec:efficiency}

We report wall-clock per-query latency on DL19 with each model,
warmup-excluded, with 100 candidates per query (Table~\ref{tab:eff}).
Aggregate cost: PAGC end-to-end pays roughly $2{\times}$ a single
pointwise pass, since RG-YN and GCCP share no LLM-side state and the
spectral-MDS step is $<1\%$ of pipeline time.

\begin{table}[t]
\centering
\small
\begin{tabular}{lcccc}
\toprule
\textbf{Model} & \textbf{RG-YN} & \textbf{GCCP} & \textbf{PAGC} & \textbf{Source} \\
\midrule
Flan-T5-Large (780M) & 2.08 & 2.39 & 4.47 & harness \\
Flan-T5-XL (3B)      & 1.99 & 2.65 & 4.64 & harness \\
Flan-UL2 (20B)       & --   & --   & 23.3 & e2e \\
\bottomrule
\end{tabular}
\caption{Mean seconds per query (100 candidates) on
NVIDIA RTX~6000~Ada, FP16, no batching. T5-Large/XL come from a
controlled latency harness with warm-up exclusion; the UL2 PAGC
number is derived from end-to-end run wall-clock (the harness run
suffered from CPU offload of UL2 weights when sharing the GPU with a
concurrent process and produced misleadingly high numbers). The
end-to-end estimate is from 21 minutes / 54 DL20 queries
($\approx 23$\,s/q end-to-end, i.e.\ RG-YN scoring $+$ GCCP scoring $+$ PAGC linear aggregation). T5-Large
and T5-XL are nearly identical in wall-clock; FP16 $+$
\texttt{device\_map='auto'} hides parameter-count differences in this
regime, but UL2 ($\sim 10\times$) is an order of magnitude slower.}
\label{tab:eff}
\end{table}

% ============================================================
\section{Discussion}
% ============================================================
\label{sec:discussion}

\paragraph{The reproduction succeeded, broadly.} Mean absolute
NDCG@10 gap of $1.6\%$ on the six TREC-DL cells and $1.9$--$4.5\%$
on BEIR (24 cells = 8 sets $\times$ 3 model scales). We exceed
the paper on five cells, including a $+3.7\%$ over-shoot on
DL20/Flan-T5-Large that earlier drafts had recorded as a deficit
(Section~\ref{sec:undocumented}). \textbf{The method's central
claim -- that an anchor-based contrastive prompt aggregated with
a relevance-grading score yields a strong pointwise reranker --
holds. What we refine is the \emph{scope} of conditions under
which the aggregation step itself remains beneficial}: the negative
DBPedia-Entity finding and the Holm-corrected non-results on E5
cells qualify the paper's framing without invalidating its central
claim.

\paragraph{The lessons are about practice, not the method.}
Eight fail-silent undocumented details between paper text and
working code is itself a finding worth communicating. Three are
make-or-break (decoder input string, target-token case, min-max
normalization): without any one, the pipeline runs to completion
but produces a broken ranking. ``Release code'' is not equivalent
to ``the paper is reproducible'': the released code is a corpus
that must be audited, and the audit is itself research labour.

\paragraph{The paper's PAGC-vs-GCCP claim is statistically more
contingent than it looks.} GCCP-alone is Holm-significantly
different from RG-YN in only 3/22 settings -- the contrastive
anchor's per-query signal becomes detectable only at high power or
high model scale. PAGC's gain over GCCP is Holm-significant in
5/22 with mixed signs: 4 positive (all BM25) and 1 negative
(DBPedia-Entity under E5). GCCP and RG-YN make complementary
mistakes mainly on noisy BM25 candidate lists; under E5 the two
agree and aggregation has little to read out. ``Aggregation is
unambiguously useful only under BM25; under E5 it is roughly a
wash and on at least one large-$n_q$ set it actively hurts''
captures the data better than the paper's framing that PAGC
reliably improves over GCCP alone.

\paragraph{First-stage retriever interacts with reranker value.}
The marginal NDCG@10 contribution of GCCP/PAGC over its first-stage
baseline depends on the baseline's quality. With BM25 the gain is
large (e.g.\ $+0.197$ PAGC over BM25 on DL20/T5-XL); with E5 it is
small ($+0.013$ on DL20) and on SciFact PAGC marginally hurts
($-0.010$). The relative value of GCCP to a first-stage upgrade
depends on the deployment regime; the original paper's BM25-only
evaluation underrepresents the strong-retriever case.

% ============================================================
\section{Conclusion}
% ============================================================

We reproduced GCCP and PAGC end-to-end, audited the implementation
against the paper, and ran several extensions and ablations the
original paper does not contain. Our findings, in order of
practical importance:

\begin{enumerate}[leftmargin=1.5em,nosep]
    \item The method's headline claim reproduces within a $1.6\%$
          mean absolute NDCG@10 gap on TREC-DL and $1.9$--$4.5\%$
          on BEIR; we exceed the paper on five cells.
    \item Eight undocumented implementation details determine whether
          a paper-only reimplementation works at all; one of them
          (target-token case) drops NDCG@10 from 0.66 to 0.24 in
          isolation, and three are make-or-break.
    \item Across 22 primary settings, paired-bootstrap (10k
          resamples) with Holm-Bonferroni correction shows:
          \textbf{GCCP-alone vs RG-YN} Holm-sig in only 3/22 (the
          contrastive-anchor signal becomes detectable only at high
          power or high model scale); \textbf{PAGC vs RG-YN}
          Holm-sig in 12/22 (all positive -- the paper's headline
          ordering holds); \textbf{PAGC vs GCCP} Holm-sig in 5/22
          with \emph{mixed signs} -- 4 positive (BM25), 1 negative
          (DBPedia-Entity, $\Delta{=}{-}0.0144$, $p_{\text{Holm}}{=}0.032$).
          Aggregation is unambiguously useful only under BM25;
          under E5 it is a wash and on at least one large-$n_q$
          set it hurts.
    \item The \textbf{spectral MDS anchor never wins} on any of
          8 BEIR sets in the paper's 4-anchor comparison and
          finishes last on 3/8. A 9-variant operationalization
          sweep identifies the paper's likely ``Top'' baseline as
          title-only Top-1 ($\Delta{-}0.003$ to paper's average);
          even after that match, Spectral is beaten by Top-3 and
          full-document Top-1, so the negative spectral conclusion
          holds.
    \item Replacing BM25 with E5 dense retrieval lifts the
          pipeline by 2--4~pts NDCG@10 on TREC-DL and exceeds both
          the paper's RG-YN+GCCP and PAGC-QSG variants on
          Flan-T5-XL. The contrastive anchor contributes less when
          the candidate list is already clean. Generalized to all
          8 BEIR sets under E5 + Flan-T5-XL, PAGC beats RG-YN
          Holm-significantly in 6/8 but \textbf{PAGC vs GCCP-alone
          is never Holm-significantly positive and is
          Holm-significantly negative on DBPedia-Entity}.
    \item Decoder-only LLMs work with our chat-template adaptation.
          \textbf{Qwen-2.5-72B-Instruct-AWQ reaches PAGC
          NDCG@10 $=0.7465$ on DL19,} surpassing our reproduced
          Flan-UL2 ($0.7095$) and both paper Flan-UL2 variants
          ($0.7206$ RG-YN+GCCP, $0.7200$ PAGC-QSG). A 50-query
          BEIR check shows the DBPedia negative \emph{strengthens}
          under this 72B backbone: PAGC ($0.4139$) underperforms
          not just GCCP-alone ($0.4380$) but also the E5
          first-stage ($0.4168$). At 7--8B, Qwen-2.5-7B is
          competitive with Flan-T5-XL ($p{=}0.125$); LLaMA-3.1-8B
          is below Flan-T5-Large; Mistral-7B underperforms.
          Backbone family matters at least as much as parameter
          count at 7--8B; quantization is not a barrier at 72B.
    \item NDCG@10 implementations are not interchangeable; a
          hand-written routine instead of \texttt{pytrec\_eval}
          accounts for $\sim 2.5$~pts of our largest BEIR gap.
\end{enumerate}

Code, per-query scores, run logs, statistical tests, and the
reproducibility checklist are released at
\url{https://github.com/utshabkg/GCCP-reproduce}.

\begin{acks}
We thank the original authors of \cite{long2025precise} for
releasing their codebase, which was essential for auditing the
implementation choices catalogued in
Section~\ref{sec:undocumented}.
We acknowledge the use of large-language-model assistants
(Anthropic Claude) for prose drafting and revision; all
experimental design, code, results, and statistical analyses
are the authors' work, and the authors take final responsibility
for all content.
\end{acks}

% ============================================================
\appendix
% ============================================================

\section{Full enumeration of the 22 primary statistical-test settings}
\label{app:full-stats}

Table~\ref{tab:full-stats} enumerates every cell in the Holm
correction family of Section~\ref{sec:stats} (auxiliary
seed-sweep and authors'-MDS-port settings re-test the same data and are
excluded). Holm uses $k=22$ within each comparison family.

\paragraph{Uncorrected counts, for calibration.}
Under raw $\alpha{=}0.05$: PAGC vs RG-YN sig in 15/22 (vs 12/22
Holm); PAGC vs GCCP in 8/22 with 1 negative (vs 5/22 Holm);
GCCP vs RG-YN in 9/22 (vs 3/22 Holm). Holm's largest effect is on
the GCCP-vs-RG-YN family (6 raw-sig drop out), consistent with
our finding that GCCP-alone's per-query signal is weaker than
aggregate-mean improvements suggest.

\begin{table*}[t]
\centering
\scriptsize
\setlength{\tabcolsep}{3pt}
\begin{tabular}{rlrrlrlrl}
\toprule
\textbf{\#} & \textbf{Setting} & \textbf{$n_q$}
& \multicolumn{2}{c}{\textbf{PAGC vs RG-YN}}
& \multicolumn{2}{c}{\textbf{PAGC vs GCCP}}
& \multicolumn{2}{c}{\textbf{GCCP vs RG-YN}} \\
\cmidrule(lr){4-5}\cmidrule(lr){6-7}\cmidrule(lr){8-9}
 & & & $\Delta$ & $p_{\text{Holm}}$ & $\Delta$ & $p_{\text{Holm}}$ & $\Delta$ & $p_{\text{Holm}}$ \\
\midrule
\multicolumn{9}{l}{\emph{TREC-DL ($n=14$)}}\\
1  & DL19 / Flan-T5-Large / BM25      & 43  & $+0.0219$ & 0.582~ns       & $+0.0512$ & $<$0.001~*** & $-0.0293$ & 1.000~ns \\
2  & DL19 / Flan-T5-XL / BM25         & 43  & $+0.0202$ & 0.582~ns       & $+0.0191$ & 0.560~ns     & $+0.0011$ & 1.000~ns \\
3  & DL19 / Flan-T5-XL / E5           & 43  & $+0.0433$ & 0.048~*        & $+0.0095$ & 1.000~ns     & $+0.0338$ & 1.000~ns \\
4  & DL19 / LLaMA-3.1-8B / BM25       & 43  & $+0.0296$ & 0.582~ns       & $+0.0202$ & 1.000~ns     & $+0.0094$ & 1.000~ns \\
5  & DL19 / Mistral-7B-v0.3 / BM25    & 43  & $+0.0713$ & 0.008~**       & $+0.0811$ & $<$0.001~*** & $-0.0098$ & 1.000~ns \\
6  & DL19 / Qwen-2.5-7B / BM25        & 43  & $+0.0685$ & $<$0.001~***   & $+0.0173$ & 0.952~ns     & $+0.0512$ & 0.144~ns \\
7  & DL19 / Qwen-2.5-72B-AWQ / BM25   & 43  & $+0.0875$ & $<$0.001~***   & $+0.0142$ & 0.784~ns     & $+0.0733$ & 0.012~*  \\
8  & DL20 / Flan-T5-Large / BM25      & 54  & $+0.0386$ & 0.132~ns       & $+0.0369$ & $<$0.001~*** & $+0.0016$ & 1.000~ns \\
9  & DL20 / Flan-T5-XL / BM25         & 54  & $+0.0213$ & 0.582~ns       & $+0.0024$ & 1.000~ns     & $+0.0189$ & 1.000~ns \\
10 & DL20 / Flan-T5-XL / E5           & 54  & $+0.0388$ & 0.318~ns       & $+0.0107$ & 1.000~ns     & $+0.0281$ & 1.000~ns \\
11 & DL20 / LLaMA-3.1-8B / BM25       & 54  & $+0.0371$ & 0.582~ns       & $+0.0301$ & 0.275~ns     & $+0.0071$ & 1.000~ns \\
12 & DL20 / Mistral-7B-v0.3 / BM25    & 54  & $+0.0623$ & 0.026~*        & $+0.0879$ & $<$0.001~*** & $-0.0255$ & 1.000~ns \\
13 & DL20 / Qwen-2.5-7B / BM25        & 54  & $+0.0281$ & 0.582~ns       & $+0.0194$ & 0.615~ns     & $+0.0087$ & 1.000~ns \\
14 & DL20 / Qwen-2.5-72B-AWQ / BM25   & 54  & $+0.0669$ & $<$0.001~***   & $-0.0003$ & 1.000~ns     & $+0.0672$ & 0.004~** \\
\midrule
\multicolumn{9}{l}{\emph{BEIR / Flan-T5-XL / E5 ($n=8$)}}\\
15 & SciFact                          & 300 & $+0.0597$ & $<$0.001~***   & $+0.0150$ & 1.000~ns     & $+0.0446$ & 0.144~ns \\
16 & NFCorpus                         & 323 & $+0.0172$ & $<$0.001~***   & $+0.0048$ & 1.000~ns     & $+0.0124$ & 0.580~ns \\
17 & TREC-COVID                       & 50  & $+0.0742$ & 0.003~**       & $+0.0107$ & 1.000~ns     & $+0.0635$ & 0.072~ns \\
18 & Touch\'e-2020                    & 49  & $+0.0477$ & 0.160~ns       & $+0.0056$ & 1.000~ns     & $+0.0421$ & 0.816~ns \\
19 & DBPedia-Entity                   & 400 & $+0.0808$ & $<$0.001~***   & $-0.0144$ & 0.032~*      & $+0.0953$ & $<$0.001~*** \\
20 & TREC-News                        & 57  & $+0.0069$ & 0.582~ns       & $+0.0024$ & 1.000~ns     & $+0.0045$ & 1.000~ns \\
21 & Robust04                         & 249 & $+0.0335$ & $<$0.001~***   & $+0.0071$ & 1.000~ns     & $+0.0264$ & 0.116~ns \\
22 & Signal-1M                        & 97  & $+0.0279$ & 0.003~**       & $-0.0025$ & 1.000~ns     & $+0.0304$ & 0.104~ns \\
\bottomrule
\end{tabular}
\caption{Full enumeration of the 22 primary settings used for the
paired-bootstrap + Holm-Bonferroni statistical tests in
Section~\ref{sec:stats}. Holm correction uses $k=22$ within each
of the three comparison families. Per-query NDCG@10 scores are
released so that any reader can recompute under their preferred
correction (e.g.\ Benjamini-Hochberg, or per-family-pair correction
across the joint $k=66$). Significance markers: \textbf{ns} =
$p_{\text{Holm}}\geq 0.05$; *: $<0.05$; **: $<0.01$; ***: $<0.001$.}
\label{tab:full-stats}
\end{table*}

% ============================================================
\bibliographystyle{ACM-Reference-Format}
\bibliography{references}

\end{document}
